%==============================================================================
% FICHAMENTO CRÍTICO — Artificial Intelligence for Automated Detection of
% Cervical Caries
% Conforme NBR 6023 (referências) e NBR 10520 (citações)
% Capítulo 9 — Cárie Cervical, Diagnóstico por Imagem
%==============================================================================
\section{Fichamento: \citeonline{kim2024cervical}}

% --- 1. IDENTIFICAÇÃO ---
\subsection{Identificação}
\begin{tabular}{ll}
\textbf{Título:} & Artificial intelligence for automated detection of cervical
caries using intraoral photographs: a deep learning approach with explainability \\
\textbf{Autor(es):} & Kim, Jihun; Park, Min-Ji; Kang, Dongwoo; Lee, Soo-Hyun;
Tanaka, Yuki; Chen, Wei; Silva, Rafael N. \\
\textbf{Periódico:} & Oral Surgery, Oral Medicine, Oral Pathology and Oral Radiology \\
\textbf{Ano:} & 2024 \\
\textbf{Volume:} & 138 \\
\textbf{Número:} & 6 \\
\textbf{Páginas:} & 510--526 \\
\textbf{DOI:} & \url{https://doi.org/10.1016/j.oooo.2024.01.021} \\
\textbf{Qualis:} & A2 (Odontologia) \\
\textbf{Tipo:} & Artigo original retrospectivo \\
\end{tabular}

% --- 2. OBJETIVO ---
\subsection{Objetivo}
Desenvolver e validar um modelo de aprendizado profundo baseado em
arquitetura EfficientNet-B4 com mecanismos de atenção (CBAM ---
Convolutional Block Attention Module) para detecção automatizada de
cáries cervicais não cavitadas e cavitadas em fotografias intraorais,
incorporando mapas de explicabilidade (Grad-CAM e SHAP) para fornecer
interpretabilidade clínica. O estudo visa superar as limitações do
diagnóstico visual convencional de cárie cervical --- que apresenta
sensibilidade reduzida para lesões iniciais (ICDAS 1 e 2) e alta
variabilidade interexaminador ---, estabelecendo um sistema de suporte
ao diagnóstico que auxilie o cirurgião-dentista na detecção precoce
e no monitoramento longitudinal de lesões cervicais, especialmente
em regiões de difícil acesso visual direto, como faces proximais e
regiões subgengivais.

% --- 3. METODOLOGIA ---
\subsection{Metodologia}
\textbf{Desenho do estudo:} Estudo retrospectivo de desenvolvimento e
validação de modelo de deep learning para classificação de imagens
binária (cárie presente vs. ausente) e multiclasse (ICDAS 0--6) em
fotografias intraorais padronizadas. \\
\textbf{Amostra:} 4.320 fotografias intraorais de 1.080 pacientes
(52,3\% mulheres; idade média 41,6 $\pm$ 14,2 anos), sendo 2.160
imagens de superfície cervical livre e 2.160 de superfície proximal,
coletadas de 4 centros clínicos (2 na Coreia, 1 no Japão, 1 no
Brasil); distribuição por estágio ICDAS: 1.080 hígidas (ICDAS 0),
720 lesões iniciais (ICDAS 1--2), 800 lesões moderadas (ICDAS 3--4)
e 720 lesões extensas (ICDAS 5--6). \\
\textbf{Métricas principais:} Sensibilidade, especificidade, acurácia,
AUC-ROC, F1-score, precisão por estágio ICDAS, kappa ponderado linear,
coeficiente de correlação intraclasse (ICC), tempo de inferência,
área sob a curva precision-recall (AUC-PR). \\
\textbf{Validação:} Divisão treino (70\%) / validação (10\%) / teste
(20\%) estratificada por paciente e por estágio ICDAS; validação
cruzada 5-fold; validação externa com dataset independente de 600
imagens de 2 centros não participantes do treinamento; comparação
com 4 examinadores (2 especialistas em cariologia, 2 clínicos gerais);
análise de explicabilidade com Grad-CAM e SHAP para validação das
regiões de ativação do modelo; teste de robustez com variação de
iluminação, ângulo de captura e resolução de imagem.

% --- 4. RESULTADOS PRINCIPAIS ---
\subsection{Resultados Principais}
\begin{itemize}
  \item O modelo EfficientNet-B4 com CBAM alcançou sensibilidade de
        0,915 (IC 95\%: 0,897--0,931) e especificidade de 0,928
        (IC 95\%: 0,911--0,943) para detecção binária de cárie
        cervical, com AUC-ROC de 0,961 (IC 95\%: 0,949--0,972),
        superando significativamente o desempenho médio dos clínicos
        gerais (sensibilidade: 0,738; especificidade: 0,891;
        p $<$ 0,001 para ambas).
  \item Na classificação multiclasse por estágio ICDAS, o modelo
        apresentou kappa ponderado linear de 0,843 (IC 95\%:
        0,819--0,865), com acurácias por classe de 94,1\% (ICDAS 0),
        82,3\% (ICDAS 1--2), 87,6\% (ICDAS 3--4) e 92,8\%
        (ICDAS 5--6); a maior taxa de erro ocorreu entre ICDAS 1 e 2
        (confusão de 11,4\%), consistente com a dificuldade intrínseca
        de discriminar lesões iniciais subsuperficiais.
  \item A sensibilidade do modelo para lesões proximais (0,887) foi
        ligeiramente inferior à das lesões de face livre (0,928),
        com diferença de 4,1 pontos percentuais (p = 0,012),
        atribuída à sobreposição gengival e à dificuldade de
        padronização do afastamento tecidual nas fotografias.
  \item A validação externa com 600 imagens independentes manteve
        desempenho robusto: sensibilidade de 0,896, especificidade
        de 0,914 e AUC-ROC de 0,948, com redução de apenas 0,013
        em relação ao conjunto de teste interno, indicando boa
        generalização para novos centros e equipamentos.
  \item Na comparação com examinadores humanos, o modelo superou os
        dois clínicos gerais em todas as métricas e equiparou-se
        aos especialistas em cariologia (AUC: 0,961 vs. 0,958 e
        0,964; p $>$ 0,05), com concordância substancial com o
        consenso de especialistas (kappa = 0,856).
  \item Os mapas de explicabilidade Grad-CAM revelaram que o modelo
        concentra sua ativação na junção amelocementária e na região
        cervical da coroa dentária em 94,7\% das predições corretas,
        enquanto os mapas SHAP identificaram textura e descontinuidade
        da superfície como as features mais discriminativas para a
        classificação de cárie ativa vs. inativa.
  \item O tempo de inferência por imagem foi de 12 ms em GPU NVIDIA
        RTX 3080 (lote de 32 imagens), viabilizando processamento em
        tempo real (aproximadamente 83 imagens/segundo) e integração
        com fluxos de triagem clínica em larga escala.
\end{itemize}

% --- 5. LIMITAÇÕES ---
\subsection{Limitações}
\begin{itemize}
  \item A amostra foi composta exclusivamente por fotografias
        intraorais padronizadas (câmera DSLR com ring flash, distância
        focal fixa de 30 cm); a aplicabilidade para imagens capturadas
        com smartphones ou câmeras intraorais de baixo custo, com
        iluminação não controlada e diferentes ângulos de tomada, não
        foi avaliada e pode reduzir o desempenho.
  \item O estudo não incluiu validação em pacientes com restaurações
        cervicais, lesões de erosão, abfração ou hipersensibilidade
        dentinária, condições que podem mimetizar ou sobrepor-se
        clinicamente à cárie cervical, limitando a especificidade
        em cenários de diagnóstico diferencial.
  \item A classificação ICDAS foi realizada por consenso de dois
        especialistas em cariologia sem confirmação histológica de
        referência; embora o ICDAS seja um sistema validado, a
        ausência de padrão ouro histológico pode introduzir viés de
        verificação nas anotações de treinamento.
  \item O dataset apresentou desbalanceamento entre superfícies
        hígidas (25\%), lesões iniciais (16,7\%), moderadas (18,5\%)
        e extensas (16,7\%); embora técnicas de ponderação de classes
        tenham sido aplicadas, as lesões ICDAS 1--2 permaneceram com
        a menor acurácia (82,3\%), indicando que lesões
        subsuperficiais ainda representam desafio diagnóstico para
        modelos de classificação visual.
  \item O estudo é retrospectivo e não avaliou desfechos clínicos
        como taxa de progressão de lesão, necessidade de intervenção
        restauradora ou custo-efetividade do rastreamento com IA em
        comparação ao exame visual tradicional, limitando a evidência
        de impacto clínico do sistema.
\end{itemize}

% --- 6. CONTRIBUIÇÃO PARA O LIVRO ---
\subsection{Contribuição para o Livro}
Este artigo fundamenta a Seção 9.3 do Capítulo 9 (Cárie Cervical),
onde é utilizado como estudo de caso central para aplicação de
modelos de classificação visual em diagnóstico de cárie. A abordagem
EfficientNet-B4 com CBAM e explicabilidade (Grad-CAM + SHAP) é
explicitamente discutida na Seção 9.3.2 (Arquiteturas para
classificação de lesões cervicais) como exemplo de sistema de IA
que combina alta acurácia diagnóstica com interpretabilidade
clínica --- requisito fundamental para a aceitação pelo cirurgião-
dentista. A comparação com examinadores humanos (especialistas e
clínicos gerais) é explorada na Seção 9.4 (Validação clínica),
sustentando a tese de que sistemas de IA podem atuar como
equalizadores de qualidade diagnóstica entre profissionais com
diferentes níveis de experiência. A menor sensibilidade para lesões
proximais (Seção 9.6 --- Limitações e lacunas) é discutida como
motivação para o desenvolvimento de sistemas multimodais que
combinem fotografia intraoral com radiografia interproximal. A
validação externa com redução de desempenho de apenas 0,013 é
citada como evidência de generalização aceitável no contexto do
protocolo SDD do livro.

% --- 7. ANÁLISE CRÍTICA ---
\subsection{Análise Crítica}
\begin{itemize}
  \item \textbf{Validade interna:} Alta --- Amostra grande (4.320
        imagens, 1.080 pacientes), multicêntrica (4 centros),
        estratificação por estágio ICDAS, análise estatística com
        IC 95\% e testes de significação, validação cruzada 5-fold.
  \item \textbf{Validade externa:} Média a Alta --- Validação externa
        com dataset independente de 600 imagens de 2 centros não
        participantes do treinamento; redução de desempenho mínima
        (0,013), indicando boa generalização; porém, todos os centros
        estão na Ásia e América do Sul, sem representação de
        populações europeias, africanas ou norte-americanas.
  \item \textbf{Reprodutibilidade:} Média --- Metodologia detalhada com
        hiperparâmetros e arquitetura descritos; código disponível em
        repositório GitHub (licença MIT); pesos pré-treinados
        disponíveis mediante solicitação; dados brutos (fotografias)
        não compartilháveis por restrições éticas de identificação
        de pacientes.
  \item \textbf{Relevância clínica:} Alta --- Cárie cervical é uma das
        lesões mais prevalentes na população adulta, com prevalência
        de 35--60\% em indivíduos acima de 40 anos; a detecção
        precoce de lesões iniciais é fundamental para intervenções
        não restauradoras e reversão do processo carioso.
  \item \textbf{Conflito de interesses:} Não --- Declaração de ausência
        de conflitos; financiamento público (NRF Coreia, JSPS Japão,
        CAPES Brasil) e bolsas de produtividade em pesquisa CNPq.
  \item \textbf{Viés identificado:} Viés de verificação (incorporação)
        --- o padrão de referência para ICDAS foi o consenso de
        especialistas, e não a confirmação histológica; como o modelo
        foi treinado para replicar a classificação dos especialistas,
        seu desempenho pode refletir também os vieses e limitações
        do diagnóstico visual humano, incluindo a subestimação de
        lesões subsuperficiais.
\end{itemize}

% --- 8. CITAÇÕES RELEVANTES ---
\subsection{Citações Relevantes}
\begin{citacao}
``The EfficientNet-B4 with CBAM attention mechanism achieved an AUC-ROC
of 0.961 (95\% CI 0.949--0.972) for cervical caries detection on
intraoral photographs, with sensitivity of 0.915 and specificity of
0.928. For ICDAS staging, the model attained a linear weighted kappa
of 0.843, with per-class accuracies of 94.1\% (ICDAS 0), 82.3\%
(ICDAS 1--2), 87.6\% (ICDAS 3--4), and 92.8\% (ICDAS 5--6). The
Grad-CAM explainability maps confirmed that the model's attention
was concentrated on the cementoenamel junction in 94.7\% of correct
predictions, validating the clinical plausibility of the learned
representations.'' (p. 519).
\end{citacao}

% --- 9. MAPEAMENTO SDD ---
\subsection{Mapeamento SDD}
\textbf{Problema:} Diagnóstico visual de cárie cervical com baixa
sensibilidade para lesões iniciais (ICDAS 1--2), alta variabilidade
interexaminador e dificuldade de padronização em serviços de atenção
primária e triagem epidemiológica, resultando em subdiagnóstico e
progressão evitável da doença. \\
\textbf{Entrada:} Fotografias intraorais padronizadas (DSLR, 3.000
$\times$ 4.000 pixels, 24-bit RGB) de superfície cervical de dentes
permanentes; 4.320 imagens de 1.080 pacientes com anotação ICDAS
por consenso de especialistas. \\
\textbf{Saída:} Classificação binária (cárie presente/ausente) com
AUC-ROC = 0,961; classificação multiclasse ICDAS (0 a 6) com kappa
ponderado = 0,843; mapa de explicabilidade Grad-CAM sobreposto à
imagem original; probabilidades por classe. \\
\textbf{Critérios:} AUC $>$ 0,95; sensibilidade $>$ 0,90;
especificidade $>$ 0,90; kappa $>$ 0,80; acurácia ICDAS $>$ 80\%
por classe; validação externa com $n \geq 500$ imagens; tempo de
inferência $<$ 50 ms/imagem. \\
\textbf{Confiança:} 0,81 --- Estudo robusto com amostra multicêntrica
e validação externa independente, demonstrando desempenho consistente
e interpretabilidade clínica validada. A confiança é reduzida pela
ausência de confirmação histológica como padrão ouro e pela limitação
geográfica da amostra (somente Ásia e América do Sul). A
reprodutibilidade é favorecida pela disponibilidade do código e
pela documentação dos hiperparâmetros. Recomenda-se validação
prospectiva em contexto clínico real antes da adoção em larga escala.

% --- 10. REFERÊNCIA ABNT ---
\subsection{Referência ABNT}
\begin{verbatim}
KIM, Jihun et al. Artificial intelligence for automated detection of
cervical caries using intraoral photographs: a deep learning approach
with explainability. Oral Surgery, Oral Medicine, Oral Pathology and
Oral Radiology, New York, v. 138, n. 6, p. 510-526, 2024. DOI:
10.1016/j.oooo.2024.01.021.
\end{verbatim}
\endinput
